\section*{Chapter \ref{ch:s2:gamma-scalping-and-event-volatility} --- Gamma Scalping and Event Volatility}

\begin{solution}{exo:s2:gamma-scalping-and-event-volatility:1}
$\sqrt{2/\pi}\times\sqrt{0.0025} = \sqrt{2/\pi}\times 0.05 = 3.99\%$.
\end{solution}

\begin{solution}{exo:s2:gamma-scalping-and-event-volatility:2}
$\sigma^2 = 0.09 + 0.0025/(10/252)$, so $\sigma = 39.1\%$.
\end{solution}

\begin{solution}{exo:s2:gamma-scalping-and-event-volatility:3}
The announcement comes outside market hours and the price gaps at the open: there is no trading between the last hedge before the report and the
first price after it, so the whole move is taken at the delta held overnight.
\end{solution}

\begin{solution}{exo:s2:gamma-scalping-and-event-volatility:4}
Most straddles pay the variance premium on diffusive days, which sellers are paid to bear. Around announcements, sellers underestimate the event's
uncertainty, more so for small, volatile, costly-to-trade firms; the event premium runs the other way for a few days.
\end{solution}

\begin{solution}{exo:s2:gamma-scalping-and-event-volatility:5}
It hedges when the delta has moved, which is when the path has delivered variance worth banking, and not on quiet intervals when a hedge trades for
little benefit.
\end{solution}

\begin{solution}{exo:s2:gamma-scalping-and-event-volatility:6}
The extra event move is a standard deviation of about 0.41\%, against an index's daily moves of about 0.7\%; an earnings gap on a stock is often
several times its daily move. The event variance priced in short index options is small and the trade's margin is correspondingly thin.
\end{solution}

\begin{solution}{exo:s2:gamma-scalping-and-event-volatility:7}
Right on average, the daily-hedged straddle loses 2.2\% per event; 15\% too low, 0.1\%; 30\% too low, it earns 2.1\%. Break-even is a little
beyond 15\% too low. The answer is not zero because the straddle pays three days of theta at a vol that carries the diffusive premium, the round-trip
option cost and the hedges before the event pays anything.
\end{solution}

\begin{solution}{exo:s2:gamma-scalping-and-event-volatility:8}
The 1.04 assumes 200 independent events a year. Earnings seasons bunch reports in a few weeks and events share market and sector moves, so the
effective number of independent bets is smaller and the Sharpe ratio lower; and one event's standard deviation is 29\% of the premium. Size on a
stress of a bad season, not on the Sharpe ratio.
\end{solution}

\begin{solution}{pb:s2:gamma-scalping-and-event-volatility:1}
\begin{enumerate}
\item $\tfrac12\Gamma S^2(\sigma_r^2 - \sigma_i^2)\,\mathrm{d}t$: gamma times the variance delivered less the variance paid for through theta.
\item The expectation is set by implied against realised variance; hedging only changes how the path's variance is collected, hence the noise.
\item $-12.1\%$ unhedged to $-14.1\%$ hourly, with standard deviations from 62.4\% down to 6.7\% and hedging costs up to 2.1\%.
\item A band rule hedges when the delta has moved by more than a threshold; a clock rule at fixed times. Bands give less noise for the cost.
\item Two expiries after the event: $\sigma^2(\tau)\tau = \sigma_d^2\tau + E$, two equations in $\sigma_d$ and $E$.
\item The \omterm{def:s2:gamma-scalping-and-event-volatility:move}{implied move} is the expected absolute event return, $\sqrt{2/\pi}\sqrt{E}$; the \omterm{def:s2:gamma-scalping-and-event-volatility:straddle}{event straddle} is bought before the event and sold after.
\item The anticipated uncertainty is large, time-varying and informative (Dubinsky and co-authors); \omterm{def:s2:gamma-scalping-and-event-volatility:straddle}{event straddles} earned 3.34\% on average from three
days before to the announcement date (Gao, Xing and Zhang).
\item Implied event variance is the true one times 0.70 times a lognormal error: the \omterm{def:s2:gamma-scalping-and-event-volatility:move}{implied move} is too low on average and wrong by a varying
amount.
\item The fall in implied vol once the event is past; a median 36.6\% to 32.0\%, 8.6\%.
\item Mean absolute return 0.83\% against 0.71\%; VIX9D fell on 65.6\% of statement days, by 3.05\% on average.
\item The event's variance is a smaller share of 30 days' variance than of 9 days'.
\item The vega lost on the event day, whatever the stock does.
\item \textbf{Named result.} The \omterm{def:s2:gamma-scalping-and-event-volatility:straddle}{event straddle} earns 2.1\% of its premium per event with a standard deviation of 29\%, from $+8.0\%$ in the fifth of
events most underpriced to $-4.5\%$ in the fifth priced about right; daily hedging maximises the return after costs (1.04 a year for 200 events).
\item 8.0, 2.7, 3.9, 0.5 and $-4.5\%$ from the most underpriced fifth to the least.
\item The delivered move is one draw from its distribution.
\item About 15\% too low in variance just to break even; 30\% for 2.1\%.
\item Independent events.
\item Quoted option prices at entry and exit, announcement dates as known then, the crush in the exit, costs, seasons that bunch events.
\item The post-event volatility sale.
\item The chance that the event moves the stock more than the market expects.
\end{enumerate}
\end{solution}

\begin{solution}{iq:s2:gamma-scalping-and-event-volatility:1}
Owning options and delta-hedging them, so that the hedges bank the path's variance while theta pays for the implied variance; it makes money when
realised volatility exceeds implied by more than the hedging costs.
\end{solution}

\begin{solution}{iq:s2:gamma-scalping-and-event-volatility:2}
From the at-the-money straddle expiring just after the report, as a share of the stock price; or, better, extract the event variance from two expiries
after the report and take $\sqrt{2/\pi}\sqrt{E}$.
\end{solution}

\begin{solution}{iq:s2:gamma-scalping-and-event-volatility:3}
Compare \omterm{def:s2:gamma-scalping-and-event-volatility:move}{implied moves} with realised event-day moves across many reports, controlling for size and volatility; trade \omterm{def:s2:gamma-scalping-and-event-volatility:straddle}{event straddles} at quoted prices
with costs; check stability across years and firm types.
\end{solution}

\begin{solution}{iq:s2:gamma-scalping-and-event-volatility:4}
Gaps larger than implied on many names at once (a season of surprises), sector-wide moves, and correlated gaps; stress with the largest past seasons
and a joint gap across a sector.
\end{solution}

\begin{solution}{iq:s2:gamma-scalping-and-event-volatility:5}
Live positions and deltas, prices, costs and the band per book; at the open prices gap and the first quotes are wide, so a hedge on the first print can
trade at a poor price; the engine needs rules for auctions and stale quotes.
\end{solution}

\begin{solution}{iq:s2:gamma-scalping-and-event-volatility:6}
$\sigma_d^2 = (\sigma_2^2\tau_2 - \sigma_1^2\tau_1)/(\tau_2 - \tau_1)$ and $E = \sigma_1^2\tau_1 - \sigma_d^2\tau_1$. An expiry before the event
carries no event variance, so the formula would treat its lower vol as diffusive and the extracted $E$ would be wrong.
\end{solution}
